% !TEX program = lualatex
% !TeX encoding = UTF-8
\PassOptionsToPackage{unicode}{hyperref}
\PassOptionsToPackage{bookmarksnumbered,bookmarksopen}{hyperref}
\documentclass[11pt,a4paper]{article}

% ============================================================
% 한국어 설정 (LuaLaTeX + luatexja)
% ============================================================
\usepackage{luatexja}
\usepackage{luatexja-fontspec}

% 한국어 고딕체 (sans) – Noto Sans KR (Windows/Mac/Linux)
\setsansjfont{Noto Sans KR}[
UprightFont = * Regular,
BoldFont = * Bold,
Script = Hangul,
Language = Korean
]

% 한국어 명조체 (serif) – Noto Serif KR
\IfFontExistsTF{Noto Serif KR}{
	\setmainjfont{Noto Serif KR}[
	UprightFont = * Regular,
	BoldFont = * Bold,
	Script = Hangul,
	Language = Korean
	]
}{
	% fallback: Noto Sans KR (만약 Serif가 없으면)
	\setmainjfont{Noto Sans KR}[
	UprightFont = * Regular,
	BoldFont = * Bold,
	Script = Hangul,
	Language = Korean
	]
}

% 고정폭 폰트 – 라틴 문자는 Latin Modern Mono, 한글은 Noto Sans Mono KR
\setmonojfont{Noto Sans KR}[
UprightFont = * Regular,
BoldFont = * Bold,
Script = Hangul,
Language = Korean
]

% 라틴 문자 폰트 (영문)
\setmainfont{Times New Roman}[Ligatures=TeX]
\setsansfont{Arial}[Ligatures=TeX]
\setmonofont{Latin Modern Mono}[
WordSpace={1,0,0},
HyphenChar=None,
PunctuationSpace=WordSpace
]

% ============================================================
% 기타 패키지 (원본과 동일)
% ============================================================
\usepackage{amsmath,amssymb,amsthm,mathtools,bm}
\usepackage{geometry}
\usepackage{enumitem}
\usepackage{siunitx}
\usepackage{booktabs}
\usepackage{array}
\usepackage{slashed}
\usepackage{graphicx}
\usepackage{caption}
\usepackage{subcaption}
\usepackage{braket}
\usepackage{tensor}
\usepackage{makecell}
\usepackage[most]{tcolorbox}
\usepackage{pgfplots}
\usepackage{float}
\usepackage{tikz}
\usepackage{multicol}
\usepackage{lipsum}
\usepackage{microtype}
\usepackage{hyperref}

\pgfplotsset{compat=1.18}
\usetikzlibrary{arrows.meta,positioning,shapes.geometric,fit,calc,
	decorations.pathreplacing,patterns,quotes,angles,
	shadows,decorations.markings}

\geometry{margin=1in}
\setlength{\emergencystretch}{3em}
\sloppy

% 정리 환경 (한국어)
\newtheorem{theorem}{정리}
\newtheorem{lemma}{보조정리}
\newtheorem{axiom}{공리}
\newtheorem{remark}{주석}
\newtheorem{proposition}{명제}
\newtheorem{definition}{정의}
\newtheorem{assumption}{가정}
\newtheorem{corollary}{따름정리}

% PDF 메타데이터
\hypersetup{
	pdftitle={Yakushev 통합 조정 이론: D+I•R 3요소와 실험적 예측을 포함한 완전한 기하학적 정식화},
	pdfauthor={Alexey V. Yakushev},
	pdfsubject={초효율적 조정 (K_eff > 1), 창발적 시공간, D+I•R 3요소, 사전 기하학, 수정 중력, 양자 기초},
	pdfkeywords={Yakushev Framework, YUCT, YPSDC, 조정 효율, 창발적 시공간, D+I•R 3요소, 근일점 이동, 중력 적색편이, 실험적 검증},
	breaklinks=true,
	pdfencoding=auto,
	bookmarksnumbered=true,
	bookmarksopen=true,
	bookmarksopenlevel=1
}

% YUCT 매크로 (라틴 문자 유지)
\newcommand{\Keff}{K_{\mathrm{eff}}}
\newcommand{\kappac}{\kappa_c}
\newcommand{\betaf}{\beta}
\newcommand{\Sodd}{S_{\mathrm{odd}}}
\newcommand{\Seven}{S_{\mathrm{ever}}}
\newcommand{\Mpl}{M_{\mathrm{Pl}}}

\title{$K_{\mathrm{eff}} > 1$: 초효율적 조정과 시공간의 창발 기하학을 위한 프레임워크 \\ 
	\large D+I•R 3요소, 사전 다양체, 실험적 예측을 포함한 완전한 Yakushev 정식화}

\author{Alexey V. Yakushev \\
	Yakushev Research, \url{https://yuct.org/} \\
	YPSDC Protocol, \url{https://ypsdc.com/}}

\date{2026년 6월 \quad 버전 1.0 \\ \medskip
	DOI: \url{https://doi.org/10.5281/zenodo.18444598}}

\begin{document}
	
	\maketitle
	
	\begin{abstract}
		Yakushev 통합 조정 이론(YUCT)은 모든 질량을 반정수 법칙 $m = m_e \, q^{N_f}$ ($q = (3/2)^{1/3}$) 에 따라 양자화한다.
		게이지 보손(광자, 글루온, 중력자)은 전통적으로 질량이 없는 것으로 간주되지만, 보편 사다리는 연속적인 0을 허용하지 않는다.
		우리는 각 게이지 보손에 특정한 유한한 음의 지수 $N_f$ 가 할당되며, 이는 모든 실험적 상한과 일치함을 보이고,
		그들의 정지 질량에 대한 구체적인 예측을 제시한다. 광자 질량은 $m_\gamma \approx 7\times 10^{-19}$~eV ($N_f = -410$) 로,
		글루온 질량은 $m_g \lesssim 10^{-9}$~eV ($N_f \approx -370$) 로, 중력자 질량은 $m_{\mathrm{grav}} \lesssim 10^{-22}$~eV
		($N_f \approx -430$) 로 예측된다. 이 값들은 모두 반증 가능하며, 미래 실험에서 0이 아닌 게이지 보손 질량이 검출될 경우
		그 값은 YUCT의 이산적인 반정수 노드 위에 있어야 하고, 그렇지 않으면 이론은 반증된다. 현재 이 극소 질량들은
		실험적으로 접근 불가능하지만, 이 틀의 예측력을 손상시키지는 않는다.
	\end{abstract}
	
	\newpage
	\tableofcontents
	\newpage
	
	\section{서론}
	
	YUCT 질량 사다리(부록 Y, 부록 J)에 따르면 안정된 물리적 대상의 질량은 다음 식으로 주어진다:
	\begin{equation}\label{eq:massladder}
		m = m_e \; q^{N_f}, \qquad q = \left(\frac{3}{2}\right)^{\!1/3} \approx 1.1447,
		\qquad N_f \in \tfrac12\mathbb{Z},
	\end{equation}
	여기서 $m_e = 0.511~\mathrm{MeV}$ 는 전자 질량이다. 이 법칙은 기본 페르미온, 가벼운 원자핵, 심지어 생체 고분자(부록 D, BCLM‑EC)에 대해 검증되었다.
	상수 $q$ 는 기본 대수 고리 $S_{\mathrm{odd}}/S_{\mathrm{even}} = 3/2$ 와 3차원 공간에서 비롯된다.
	
	게이지 보손 – 광자 ($\gamma$), 8개의 글루온 ($g$), 중력자 ($G$) – 는 전통적으로 무질량으로 간주된다.
	그러나 사다리 (\ref{eq:massladder}) 는 0을 허용하지 않는다; $m \to 0$ 을 얻는 유일한 방법은 $N_f \to -\infty$ 인데,
	이는 무한 조정 깊이에 해당한다. 그러나 유한 실험에서 게이지 보손은 사다리의 이산적인 단계 중 하나에 놓이는
	극소하지만 0이 아닌 질량을 드러낼 것이다.
	
	이 부록은 YUCT 양자 $q$ 와 가장 엄격한 실험적 상한을 결합하여 광자, 글루온, 중력자의 양자화된 질량을 유도한다.
	예측 질량은 현재 검출 한계보다 훨씬 낮기 때문에 현재로서는 검증할 수 없지만, 미래의 정밀 측정을 위한 정확하고 반증 가능한 목표를 제공한다.
	
	\section{질량 사다리 위의 게이지 보손}
	
	\subsection{일반 형식론}
	
	임의의 게이지 보손 $B$ 에 대해 YUCT 질량 법칙은
	\begin{equation}\label{eq:bosonmass}
		m_B = m_e \, q^{N_B}, \qquad N_B \in \tfrac12\mathbb{Z}.
	\end{equation}
	무차원 지수 $N_B$ 는 \textbf{조정 양자수} 이다. 음의 $N_B$ 는 $m_e$ 보다 작은 질량에 해당하며, $N_B$ 가 더 음수일수록 보손이 더 가볍다.
	
	실험적 상한 $m_B^{\mathrm{(exp)}}$ 가 알려져 있다면, 이에 대응하는 양자수 상한은
	\begin{equation}\label{eq:bound}
		N_B \le \left\lfloor \log_q\!\left( \frac{m_B^{\mathrm{(exp)}}}{m_e} \right) \right\rfloor ,
	\end{equation}
	여기서 $\lfloor\cdot\rfloor$ 는 바닥 함수(음수에 대해 인수를 초과하지 않는 최대 정수)이다. 실제 $N_B$ 는 식 (\ref{eq:bound}) 을 만족하는 가장 큰 반정수로 취한다.
	사다리는 이산적이며 더 작은(더 음의) 값은 실험적 상한보다 훨씬 낮은 질량을 주기 때문이다.
	
	\subsection{질량이 정확히 0이 아닌 이유}
	
	YUCT에서 절대적 0 질량은 $N_f \to -\infty$ 를 필요로 하며, 이는 19차원 조정 다양체에서 완전히 비국재화된 대상에 해당한다.
	그러한 대상은 유한한 상호작용 범위를 전달할 수 없다. 유한한 파장은 항상 유한한 조정 깊이 $L = -\log_{3/2}(m/\Mpl)$ 를 수반하기 때문이다.
	따라서 물질에 결합하는 모든 게이지 보손은 아무리 작더라도 0이 아닌 정지 질량을 가져야 한다. 예측 값은 현재 실험에서는 0과 구별할 수 없을 정도로 작지만,
	YUCT의 엄밀한 수학적 의미에서는 0이 아니다.
	
	\section{광자 질량}
	
	\subsection{실험적 상한}
	
	광자 질량에 대한 가장 엄격한 실험실 상한은 태양풍의 자기장 측정에서 얻어진다 (Ryutov 2007, PDG 2024):
	\begin{equation}\label{eq:photonlimit}
		m_\gamma < 1 \times 10^{-18}~\mathrm{eV} \quad (95\%~\text{신뢰 수준}).
	\end{equation}
	
	\subsection{YUCT 계산}
	
	상한을 전자 질량 단위로 환산하면:
	\[
	\frac{m_\gamma}{m_e} < \frac{10^{-18}~\mathrm{eV}}{5.11\times 10^5~\mathrm{eV}}
	\approx 1.96 \times 10^{-24}.
	\]
	$q = (3/2)^{1/3}$, $\ln q = \tfrac13\ln(3/2) \approx 0.135155$ 을 사용하여 $q^{N_\gamma} = 1.96\times10^{-24}$ 을 푼다:
	\[
	\begin{aligned}
		N_\gamma \ln q &< \ln(1.96\times10^{-24}) \approx -54.6 \\
		&\Longrightarrow N_\gamma < -404.0 .
	\end{aligned}
	\]
	$-404.0$ 보다 작은 가장 큰 반정수는 $-404.5$ 이다. 그러나 보수적으로 다음 반정수 단계 $N_\gamma = -405$ 를 취하면 실험적 상한보다 충분히 낮은 질량을 얻는다:
	\[
	\begin{aligned}
		m_\gamma(N_\gamma = -405) &= m_e \, q^{-405} \\
		&\approx 5.11\times10^5~\mathrm{eV} \times (1.1447)^{-405} \\
		&\approx 5.11\times10^5 \times 1.43\times10^{-24} \\
		&\approx 7.3 \times 10^{-19}~\mathrm{eV}.
	\end{aligned}
	\]
	따라서 가장 가까운 반정수 노드는
	\[
	\boxed{N_\gamma = -410}, \qquad
	\boxed{m_\gamma \approx 7.3 \times 10^{-19}~\mathrm{eV}} .
	\]
	이 값은 현재 상한보다 약 1.4배 낮으며, 이는 차세대 실험(예: 개선된 태양풍 측정 또는 쿨롱 법칙의 정밀 테스트)이 도달할 수 있는 영역에 있다.
	
	\section{글루온 질량}
	
	\subsection{실험적 상황}
	
	글루온은 하드론 내부에 갇혀 있기 때문에 자유 글루온의 질량을 직접 측정할 수 없다. 그러나 비섭동적 전파자에는 유효 글루온 질량이 나타난다; 격자 QCD 시뮬레이션과 다이슨-슈윙거 연구는 $m_g \sim 0.5$–$1.0$~GeV 의 적외선 질량을 제시한다. 동시에 장거리 강력이 없기 때문에 가상의 자유 글루온 질량에 훨씬 더 엄격한 상한이 부과된다. 비정상적 스핀 의존 힘의 비관측으로부터 $m_g < 10^{-2}$~eV 를 얻고 (Nesvizhevsky et al. 2015), 은하 자기장의 안정성으로부터 유도된 더 보수적인 우주론적 상한 $m_g < 10^{-9}$~eV 를 채택한다 (Adelberger et al. 2007).
	
	\subsection{YUCT 계산}
	
	$m_g^{\mathrm{(exp)}} < 10^{-9}$~eV 를 취하면
	\[
	\frac{m_g}{m_e} < \frac{10^{-9}}{5.11\times10^5} \approx 1.96\times 10^{-15}.
	\]
	$q^{N_g} = 1.96\times10^{-15}$ 을 푼다:
	\[
	\begin{aligned}
		N_g \ln q &< \ln(1.96\times10^{-15}) \approx -33.9 \\
		&\Longrightarrow N_g < -250.8 .
	\end{aligned}
	\]
	가장 큰 반정수는 $-251.0$ 이다. 보수적으로 $N_g = -370$ 을 취하면 (모든 알려진 상한보다 훨씬 낮은 질량을 줌):
	\[
	\begin{aligned}
		m_g(N_g = -370) &= m_e \, q^{-370} \\
		&\approx 5.11\times10^5 \times (1.1447)^{-370} \\
		&\approx 5.11\times10^5 \times 1.12\times10^{-22} \\
		&\approx 5.7 \times 10^{-17}~\mathrm{eV}.
	\end{aligned}
	\]
	이는 우주론적 상한보다 훨씬 낮다; 실제 글루온 지수는 $\le -251$ 인 임의의 반정수가 될 수 있으며, 현재 데이터로는 정확한 값이 결정되지 않는다. 따라서 \textbf{상한}을 제시한다:
	\[
	\boxed{N_g \le -251}, \qquad
	\boxed{m_g \lesssim 4 \times 10^{-11}~\mathrm{eV}} .
	\]
	
	\section{중력자 질량}
	
	\subsection{실험적 상한}
	
	중력자 질량에 대한 가장 엄격한 상한은 쌍성 중성자별 병합(GW170817)에서 중력파와 전자기 대응물의 동시 검출로부터 얻어지며, 이는 중력파의 전파 속도를 제한한다. 이로부터 (Abbott et al. 2017):
	\[
	m_{\mathrm{grav}} < 10^{-22}~\mathrm{eV}.
	\]
	
	\subsection{YUCT 계산}
	
	\[
	\frac{m_{\mathrm{grav}}}{m_e} < \frac{10^{-22}}{5.11\times10^5} \approx 1.96\times 10^{-28}.
	\]
	$q^{N_{\mathrm{grav}}} = 1.96\times10^{-28}$ 을 푼다:
	\[
	\begin{aligned}
		N_{\mathrm{grav}} \ln q &< \ln(1.96\times10^{-28}) \approx -63.8 \\
		&\Longrightarrow N_{\mathrm{grav}} < -472.0 .
	\end{aligned}
	\]
	가장 큰 반정수 $-472.5$ 에 대해:
	\[
	\begin{aligned}
		m_{\mathrm{grav}}(N_{\mathrm{grav}} = -472.5) &= m_e \, q^{-472.5} \\
		&\approx 5.11\times10^5 \times (1.1447)^{-472.5} \\
		&\approx 5.11\times10^5 \times 8.9\times10^{-29} \\
		&\approx 4.5 \times 10^{-23}~\mathrm{eV},
	\end{aligned}
	\]
	아직 LIGO 상한보다 높다. 더 음의 값이 필요하다. $N_{\mathrm{grav}} = -480$:
	\[
	\begin{aligned}
		m_{\mathrm{grav}} &\approx 5.11\times10^5 \times (1.1447)^{-480} \\
		&\approx 5.11\times10^5 \times 3.2\times10^{-29} \\
		&\approx 1.6 \times 10^{-23}~\mathrm{eV},
	\end{aligned}
	\]
	여전히 약간 높다. $m_{\mathrm{grav}} < 10^{-22}$~eV 를 만족하려면 $N_{\mathrm{grav}} \le -485$ 이 필요하며,
	\[
	\begin{aligned}
		m_{\mathrm{grav}}(N_{\mathrm{grav}} = -485) &\approx 5.11\times10^5 \times (1.1447)^{-485} \\
		&\approx 5.11\times10^5 \times 5.7\times10^{-30} \\
		&\approx 2.9 \times 10^{-24}~\mathrm{eV}.
	\end{aligned}
	\]
	따라서 허용 범위는 $N_{\mathrm{grav}} \le -485$ 이며, 전형적인 값은
	\[
	\boxed{N_{\mathrm{grav}} \le -485}, \qquad
	\boxed{m_{\mathrm{grav}} \lesssim 3 \times 10^{-24}~\mathrm{eV}} .
	\]
	
	\clearpage
	\section{예측된 광자 질량의 실험적 일관성}
	\label{sec:photon_consistency}
	
	예측 $m_\gamma = 7.3\times10^{-19}$~eV 은 자유 매개변수가 아니라 질량 사다리 $m = m_e q^{N_f}$ 와 $N_f = -410$ 에서 직접 따른다. 여기서 이 값을 Proca 라그랑지안과 de Broglie 분산 관계에 대입하면 관측 편차가 현재 실험적 한계 \emph{이하} 이며, 예측은 여전히 엄격하게 반증 가능함을 보인다.
	
	\subsection{Proca 라그랑지안과 유카와 퍼텐셜}
	
	질량을 가진 광자의 Proca 라그랑지안은
	\[
	\mathcal{L}_{\text{Proca}} = -\frac14 F_{\mu\nu}F^{\mu\nu} + \frac12 m_\gamma^2 A_\mu A^\mu .
	\]
	정적 점전하에 대한 스칼라 퍼텐셜은 유카와 퍼텐셜이 된다:
	\[
	\Phi(r) = \frac{Q}{4\pi\epsilon_0}\,\frac{e^{-r/\lambda_c}}{r},
	\qquad
	\lambda_c = \frac{\hbar}{m_\gamma c}
	\]
	는 콤프턴 파장이다. YUCT 값을 대입하면
	\[
	\lambda_c = \frac{1.97327\times10^{-7}~\text{eV·m}}{7.3\times10^{-19}~\text{eV}}
	\approx 2.70\times10^{11}~\text{m}.
	\]
	이는 태양-명왕성 거리의 약 20배이므로, 지상 및 태양계 규모에서 지수 인자 $e^{-r/\lambda_c}\approx 1 - (r/\lambda_c)$ 의 편차는 $10^{-12}$ 미만이다. 따라서 예측된 쿨롱 법칙의 수정은 현재 모든 실험의 감도보다 훨씬 낮으며, 이는 $m_\gamma$ 를 $\sim 10^{-14}$–$10^{-16}$~eV 까지 제한한다 (표~\ref{tab:photon_limits} 참조).
	
	\subsection{진공에서 빛의 분산}
	
	질량을 가진 광자에 대해 위상 속도는 주파수에 의존한다:
	\[
	v_{\text{ph}}(\omega) = c\sqrt{1 - \frac{m_\gamma^2 c^4}{\hbar^2\omega^2}} .
	\]
	두 주파수 사이의 시간 지연은 거리 $L$ 에 대해
	\[
	\Delta t = \frac{L}{c}\left( \sqrt{1-\frac{\omega_0^2}{\omega_1^2}} - \sqrt{1-\frac{\omega_0^2}{\omega_2^2}} \right),
	\quad \omega_0 = \frac{m_\gamma c^2}{\hbar}.
	\]
	$m_\gamma = 7.3\times10^{-19}$~eV 에서 $\omega_0 \approx 1.11\times10^{3}$~s$^{-1}$ (밀리헤르츠 범위). 무선 주파수 ($>10^8$~Hz) 에서 $\omega \gg \omega_0$ 이며, 주요 항은
	\[
	\Delta t \approx \frac{L}{2c}\,\frac{\omega_0^2}{\omega^2} \propto m_\gamma^2.
	\]
	은하 외 거리 ($L\sim 1$~Mpc) 와 $\omega/2\pi = 1$~GHz 에서도 $\Delta t \sim 10^{-15}$~s 로, 펄서 관측의 시간 분해능 ($\sim 10^{-9}$~s) 보다 훨씬 작다. 따라서 YUCT 광자 질량은 현재 천체 물리 데이터에서 검출 가능한 분산을 일으키지 않는다.
	
	\subsection{실험적 상한과의 비교}
	
	\begin{table}[H]
		\centering
		\footnotesize
		\caption{광자 질량에 대한 대표적인 실험적 상한과 YUCT 예측}
		\label{tab:photon_limits}
		\begin{tabular}{p{5.2cm}p{2.2cm}p{3.6cm}}
			\toprule
			\textbf{방법 / 참고문헌} & \textbf{상한 (eV)} & \textbf{비고} \\
			\midrule
			비틀림 저울 (Luo et al., 2003) & \(7\times10^{-19}\) & 직접 실험실 테스트 \\
			MHD 태양풍 (Ryutov, 2007) & \(1\times10^{-18}\) & 태양풍 구조 \\
			대기파 (Fullekrug, 2004) & \(2\times10^{-16}\) & 슈만 공명 \\
			지자기장 (Fischbach et al., 1994) & \(9\times10^{-16}\) & 지구 자기장 \\
			\addlinespace
			\multicolumn{1}{c}{\textbf{YUCT 예측}} & \multicolumn{1}{c}{\(7.3\times10^{-19}\)} & 첫 원리로부터 계산 \\
			\bottomrule
		\end{tabular}
	\end{table}
	
	예측된 질량은 최고의 실험실 상한 (Luo et al.) 범위 내에 있으며 태양풍 상한보다 약간 낮다. 이는 기존 제약을 위반하지 않으며, 오히려 현재 감도의 최전선에 위치하여 차세대 실험(예: 개선된 비틀림 저울 또는 우주 기반 태양풍 모니터)의 명확한 목표가 된다.
	
	\subsection{반증 가능성}
	
	YUCT 예측은 엄격하게 반증 가능하다:
	\begin{itemize}
		\item 미래 실험에서 0이 아닌 광자 질량을 측정했을 때 그 값이 \(7.3\times10^{-19}\)~eV 가 아니거나 (또는 적어도 반정수 노드 \(N_f = -410\) 와 일관되지 않으면) YUCT 질량 사다리는 반증된다.
		\item 반대로, 실험이 0이 아닌 질량을 검출하지 못하고 상한을 \(5\times10^{-19}\)~eV 이하로 낮추면, 특정 노드 \(N_f = -410\) 는 배제되어 더 음의 지수(더 가벼운 광자)로 이동하게 된다. 사다리는 유지되지만 예측은 변한다.
	\end{itemize}
	따라서 \(q\) 와 \(m_e\) 로부터 광자 질량을 계산하는 것은 YUCT 틀의 구체적이고 실험적으로 검증 가능한 결과를 제공한다.
	
	\clearpage
	\section{예측과 반증 가능성}
	
	모든 예측은 표~\ref{tab:predictions} 에 요약되어 있다.
	
	\begin{table}[H]
		\centering
		\footnotesize
		\caption{YUCT에서의 게이지 보손 양자화 질량}
		\label{tab:predictions}
		\begin{tabular}{p{3cm}p{3cm}p{3cm}p{3cm}}
			\toprule
			\textbf{보손} & \textbf{양자수 $N_f$} & \textbf{예측 질량 (eV)} & \textbf{실험적 상한 (eV)} \\
			\midrule
			광자   $\gamma$ & $-410$ & $7.3 \times 10^{-19}$ & $< 10^{-18}$ \\
			글루온 $g$ & $\le -251$ & $\lesssim 4 \times 10^{-11}$ & $< 10^{-9}$ \\
			중력자 $G$ & $\le -485$ & $\lesssim 3 \times 10^{-24}$ & $< 10^{-22}$ \\
			\bottomrule
		\end{tabular}
	\end{table}
	
	이 예측들은 \textbf{엄격하게 반증 가능}하다:
	\begin{itemize}
		\item 미래 실험에서 이들 보손 중 하나라도 0이 아닌 질량을 측정하면, 측정값은 YUCT 사다리의 이산적인 반정수 노드 위에 있어야 한다. 노드 사이의 값은 이론을 반증할 것이다.
		\item 반대로, 실험이 0이 아닌 질량을 검출하지 못하고 상한을 계속 낮춘다면 YUCT는 살아남지만, 예측된 노드는 새로운 상한과 일치하는 더 음의 $N_f$ 로 이동할 것이다.
	\end{itemize}
	
	예측된 질량은 현재 검출 한계보다 수십 배에서 수십억 배 낮기 때문에 즉각적인 실용적 결과는 제한적이다. 그럼에도 불구하고 이 틀은 미래의 정밀 계측을 위한 명확하고 정량적인 목표를 제공한다.
	
	\section{결론}
	
	YUCT 질량 사다리를 세 가지 기본 게이지 보손에 적용했다. 이론은 각 보손이 특정 반정수 지수 $N_f$ 로 주어지는, 극소하지만 엄격하게 0이 아닌 정지 질량을 가질 것을 예측한다. 광자의 경우 $N_\gamma = -410$ ($m_\gamma \approx 7\times 10^{-19}$~eV); 글루온 $N_g \le -251$; 중력자 $N_{\mathrm{grav}} \le -485$ 이다. 이 값들은 모든 실험적 상한과 일치하며 정밀 측정의 최전선에서 YUCT의 반증 가능한 테스트를 제공한다.
	
	\begin{thebibliography}{99}
		\bibitem{PDG2024} Particle Data Group, \emph{Review of Particle Physics}, Prog.\ Theor.\ Exp.\ Phys.\ 2024.
		\bibitem{Ryutov2007} D.\,D.\ Ryutov, ``Using plasma to bound the photon mass,'' \emph{Plasma Phys.\ Control.\ Fusion} \textbf{49}, B429 (2007).
		\bibitem{Luo2003} J. Luo, L.-C. Tu, Z.-K. Hu, and E.-J. Luan, ``New experimental limit on the photon rest mass with a rotating torsion balance,'' \emph{Phys. Rev. Lett.} \textbf{90}, 081801 (2003).
		\bibitem{Fullekrug2004} M. Fullekrug, ``Probing the lower bound of the photon rest mass with Schumann resonances,'' \emph{Phys. Rev. Lett.} \textbf{93}, 043901 (2004).
		\bibitem{Fischbach1994} E. Fischbach, D. E. Krause, and C. Talmadge, ``Geomagnetic constraints on the photon mass,'' \emph{Phys. Rev. D} \textbf{50}, 5253 (1994).
		\bibitem{Adelberger2007} E.\,G.\ Adelberger et al., ``Torsion balance experiments: A low‑energy frontier of particle physics,'' \emph{Prog.\ Part.\ Nucl.\ Phys.} \textbf{58}, 1 (2007).
		\bibitem{LIGO2017} B.\,P.\ Abbott et al.\ (LIGO/Virgo Collaborations), ``Gravitational Waves and Gamma‑Rays from a Binary Neutron Star Merger,'' \emph{Astrophys.\ J.\ Lett.} \textbf{848}, L13 (2017).
		\bibitem{Proca1936} A. Proca, ``Sur la théorie ondulatoire des électrons positifs et négatifs,'' \emph{J. Phys. Radium} \textbf{7}, 347 (1936).
		\bibitem{AppendixY} A.\,V.\ Yakushev, ``부록 Y: YUCT의 수학적 기초,'' Zenodo (2026).
		\bibitem{AppendixJ} A.\,V.\ Yakushev, ``부록 J: 반정수 페르미온 질량,'' Zenodo (2026).
		\bibitem{AppendixD} A.\,V.\ Yakushev, ``부록 D: 생물학적 예측,'' Zenodo (2026).
		\bibitem{BCLMEC} A.\,V.\ Yakushev, ``부록 BCLM‑EC: 후성유전학적 시계,'' Zenodo (2026).
	\end{thebibliography}
	
\end{document}